% --- Pacotes de Geometria e Tipografia ---
\usepackage[a4paper,lmargin=1.5cm,rmargin=1.5cm,tmargin=1.2cm,bmargin=1.5cm]{geometry}
\usepackage[brazilian]{babel}
\usepackage[T1]{fontenc}
\usepackage{microtype}

% --- Matemática e Símbolos ---
\usepackage{amsmath,amsfonts,amssymb}
\usepackage{mathtools}
\usepackage{latexsym}
\usepackage{esint}

% --- Gráficos e Diagramas ---
\usepackage{graphicx}
\usepackage{tikz}
\usetikzlibrary{calc,arrows.meta,shadings}
\usepackage{pgfplots}
\pgfplotsset{compat=1.18}

% --- Elementos Visuais e Layout ---
\usepackage[most]{tcolorbox}
\usepackage{enumitem}
\usepackage[normalem]{ulem}
\usepackage[Smaller]{cancel}
\usepackage{xcolor}
\usepackage[nodisplayskipstretch]{setspace}

% --- Navegação e Links ---
\usepackage[colorlinks=true, linkcolor=black, urlcolor=blue]{hyperref}
\usepackage{bookmark}

% --- Metadados do Documento ---
\hypersetup{
	pdftitle={Solucionário: Variedades Diferenciáveis},
	pdfauthor={Anselmo Ribeiro Lopes},
	pdfkeywords={Variedades Diferenciáveis, Geometria Diferencial, PPGMAT UFCG},
	pdfcreator={TeXstudio},
	pdfproducer={LuaLaTeX com motor LuaHBTeX},
	pdfsubject={Licenciado sob CC BY-NC-SA 4.0. Este material é um registro acadêmico da disciplina de Variedades Diferenciáveis - PPGMAT/UFCG.},
	pdfstartview={FitH},
	pdflang={pt-BR},
	hidelinks,
}

% --- Configurações Globais de Estilo ---
\setstretch{1.3}
\setlength{\parindent}{0pt}
\relpenalty=10000
\binoppenalty=10000
\everymath{\displaystyle}
\pagestyle{empty}
\renewcommand{\CancelColor}{\color{red}}
%\setlength{\overfullrule}{5pt} % Desenha uma barra de 5pt no final de hboxes

% --- Contador e Comandos Customizados do Dr. Geometros ---
\newcounter{numquestao}

% Comando \questao: Enunciado em tcolorbox cinza
% Altere a definição da \questao para isso:
\newcommand{\questao}[1]{%
	\stepcounter{numquestao}%
	% 1. Criamos um destino único (ex: L1.1, L2.1) para a navegação não se perder
	\hypertarget{dest.\listid.\arabic{numquestao}}{}%
	% 2. Criamos o bookmark apontando explicitamente para esse destino
	\bookmark[level=1, dest={dest.\listid.\arabic{numquestao}}]{Questão \arabic{numquestao}}%
	\begin{tcolorbox}[
		enhanced,
		sharp corners,
		boxrule=0.5pt,
		colback=gray!6!white, 
		colframe=black,         
		left=1pt, right=1pt, top=1pt, bottom=1pt,
		width=1\textwidth,
		before skip=3pt,
		after skip=3pt,
		fontupper=\normalfont
		]%
		\mbox{\textbf{Questão \arabic{numquestao}.}}\hspace{0.3em}\ignorespaces#1%
	\end{tcolorbox}%
}

% Comando \resultadoextra: Teoremas ou definições auxiliares
\newcommand{\resultadoextra}[2]{%
	\begin{tcolorbox}[
		enhanced,
		sharp corners,
		boxrule=0.5pt,
		colback=gray!6!white, 
		colframe=black,         
		left=1pt, right=1pt, top=1pt, bottom=1pt,
		width=1\textwidth,
		before skip=3pt,
		after skip=3pt,
		fontupper=\normalfont
		]%
		\mbox{\textbf{#1}}\hspace{0.3em}\ignorespaces#2%
	\end{tcolorbox}%
}

% Comando \solucao: Resolução padrão terminando com quadradinho
\newcommand{\solucao}[1]{%
	\par\vspace{5pt}
	\textbf{\underline{Solução}:} #1 \hfill$\blacksquare$
	\par\vspace{5pt}
}

% Comando \solucaoalt: Soluções alternativas/estruturais
\newcommand{\solucaoalt}[1]{%
	\par\vspace{5pt}
	\textbf{\underline{Solução Alternativa}}: #1 \hfill$\blacksquare$
	\par\vspace{5pt}
}

% --- Operadores Matemáticos e Comandos Rápidos ---
\newcommand{\cqd}{\hfill$\blacksquare$}
\DeclareMathOperator{\supp}{supp}
\DeclareMathOperator{\sen}{sen}
\DeclareMathOperator{\Img}{Im}
\DeclareMathOperator{\dist}{dist}
\newcommand{\subsubset}{\subset\subset}
\DeclareMathOperator{\Inv}{Inv}